\section{The joint lower measure and its transport}
\label{app:lower-details}

This appendix supplies the probability and measure constructions used in
Section~\ref{lower:chapter}.  We keep the original physical blocks
throughout.  Conditional on its uncolored total \(n_i\), a block has
\(n_i\) uniform points on its interval and all \(n_i+1\) spacings.
Its fine labels have their original categorical law.  Conditional on
a complete named count vector, they are a uniform multiset permutation.
Constants in this appendix depend on the fixed dimension and the fixed
compact scalar, count, and wall parameters.

\subsection{Revealing part of an original word}

The following elementary estimate lets us retain an entire coupled
count band when an original comparison class reveals some of its labels.

\begin{lemma}[Central counts after partial revelation]
\label{ld:reveal}
In block \(i\), let \(N_i\) labels be independent with alphabet size
\(a_i\le H\) and probabilities \(p_{i,c}\ge p_*>0\).  The blocks are
independent under their original law \(\pi\).  Let \(A(K)\ge0\) be
a finite weight on their complete count vectors, supported on
\[
 |K_{i,c}-N_ip_{i,c}|\le B\sqrt{1+N_i},
 \qquad W=\mathbb E_\pi A.
\]
Fix \(m_i\) slots in block \(i\).  For any positive-probability event
\(Q\) measurable in just these slots, including coupled or partial
signatures,
\begin{equation}
 \mathbb E_\pi[A\mid Q]\le CW
 \prod_{i:m_i>0}
 \left(\frac{1+N_i}{1+N_i-m_i}\right)^{(a_i-1)/2}
 \le CW\prod_{i:m_i>0}(1+m_i)^{(a_i-1)/2}.
 \label{ld:reveal-bound}
\end{equation}
If \(m_i\le N_i/2\) in every affected block, the ratio is bounded
by a fixed constant.  A block with no revealed slots incurs no factor.
\end{lemma}

\begin{proof}
For a fixed \(a\)-letter law in this compact simplex, its multinomial
mass \(P_{n,p}(k)\), with \(d=a-1\), satisfies
\begin{equation}
 \sup_kP_{n,p}(k)\le C(1+n)^{-d/2},\qquad
 P_{N,p}(k)\ge c(1+N)^{-d/2}
 \quad\text{on the stated central window}.
 \label{ld:multinomial-atoms}
\end{equation}
For the first bound, Fourier inversion, with the last coordinate set
to zero, bounds a coefficient by the integral of
\(\lvert p_a+\sum_{j\le d}p_je^{it_j}\rvert^n\).  The square of
this modulus is at most
\[
 1-2p_*^2\sum_{j\le d}(1-\cos t_j)
       \le1-c\sum_{j\le d}t_j^2,\qquad |t_j|\le\pi.
\]
Its integral is \(O((1+n)^{-d/2})\).  For the lower bound, Stirling's
inequalities give
\(P_{N,p}(k)\ge cN^{-d/2}\exp\{-ND(k/N\Vert p)\}\) when
\(N\) is large.  The entropy Hessian is bounded on the segment between
\(p\) and \(k/N\), so its exponent is \(O(B^2)\).  Bounded \(N\)
is covered by compactness and feasibility.  The case \(a=1\) is
deterministic.

Condition first on every revealed label, with count vector \(b_i\).
The remaining labels still have their original independent laws, and
\[
 \pi(K=k\mid\text{revealed labels})
       =\prod_iP_{N_i-m_i,p_i}(k_i-b_i).
\]
For \(m_i=0\), its factor cancels exactly against the original type
mass.  For the other blocks, \eqref{ld:multinomial-atoms} bounds their
ratio by the first product in \eqref{ld:reveal-bound}.  Multiply by
\(A(k)\) and sum the whole coupled type vector.  Finally, conditioning
on \(Q\) is a convex mixture of these revealed assignments.
The last inequality follows from
\((1+N)/(1+N-m)\le1+m\).  If \(W=0\), all the relevant feasible
type weights vanish, and the conclusion is immediate.
\end{proof}

For a fixed original history and anchor, a literal signature class is
an event \(Q\) of this kind: it reads only the slots with nontrivial
partitions.  This applies to that fixed class under the original law.
It does not change that law to a law conditioned on an adaptively chosen
history.  In particular, if a class fully exposes blocks in \(I\) and
reads \(m_j\) additional slots elsewhere, its count cost is at most
\[
 C\prod_{i\in I}(1+N_i)^{(a_i-1)/2}
       \prod_{j\notin I}(1+m_j)^{(a_j-1)/2}.
\]
Only counts in actually revealed blocks or strips appear.

\subsection{Two free intervals in the original class}

We next bound the path part after the complete type and exposed labels
have been fixed.  Consider one selected block of length \(P\), total
\(n\asymp P\), and binary risk
\[
 R(t)=ct-aN(t)+bN_B(t),\qquad c,b,a-b>0,
\]
with these parameters in fixed compact intervals.  Every accepted word
has an actual endpoint pair satisfying
\begin{equation}
 0\le x\le C_{\rm box}X,\quad 0\le y\le C_{\rm box}Y,
 \quad y-x=cP-an+bk,\quad x+R(t)\ge-h,
 \label{ld:bank-path}
\end{equation}
where \(X,Y\ge h_*>0\).  All additional walls and refinements can
remain in the acceptance.  Assume the uncolored endpoint bound
\begin{equation}
 N([0,t])+N([P-t,P])\le C_{\rm occ}(1+t),
       \qquad 0\le t\le P/2.
 \label{ld:endpoint-occupancy}
\end{equation}
The original class may read prescribed slots in endpoint strips of
width \(C_s(1+d)\), but is unrestricted between them.  This last
condition is a statement about its literal signatures.

Put
\begin{equation}
 \mathcal T=\frac{(1+X)(1+Y)}n,\qquad
 p_0=\frac{a-cP/n}{b},\qquad
 \rho_0=\frac{(cP/(n+1))^2}
 {(cP/(n+1))^2+b^2p_0(1-p_0)}.
 \label{ld:bank-parameters}
\end{equation}
On the rare branch \(\mathcal T<t_0\), the accepted interior type
and \eqref{ld:bank-path} imply that \(p_0\) belongs to a fixed compact
subinterval of \((0,1)\).  Define
\(\Psi=\mathcal T^{2\chi(\eta,\rho_0)}\) on this branch and
\(\Psi=1\) otherwise, where \(\eta\) lies in a fixed compact
subinterval of \((0,1)\).

Fix \(0<\gamma<1/2\) and
\[
 D_n=[C_{\rm cut}\log(2+n)]^{1/\gamma}.
\]
Set \(q_d=1\) if the branch is saturated, if \(n<n_0\), or if
\(d>D_n\).  Here \(n_0\) is chosen so that, in the remaining cases,
the endpoint strips contain at most a fixed small fraction of \(n\)
slots and the following two intervals fit in the free middle.  Choose
their starting and ending boundaries \(s_d,t_d\) just inside those
strips.  Write
\begin{equation}
 U_d=1+X+d,\quad V_d=1+Y+d,\quad M_d=\max(U_d,V_d),\quad
 L_d=\varepsilon_0\min\{P,(n/M_d)^2\}.
 \label{ld:bank-length}
\end{equation}
Take intervals \([s_d,s_d+L_d]\), \([t_d-L_d,t_d]\).
The constant \(\varepsilon_0\) is small enough that they are disjoint
and their combined occupancy is at most half the unread slots, by
\eqref{ld:endpoint-occupancy}.  Grid rounding changes only fixed
allowances.  If the resulting length is below a fixed scalar threshold,
again set \(q_d=1\).

In the remaining case set
\[
 p_\pm=p_0\pm C_1(X+Y+1+d)/n.
\]
Reduce \(t_0\) and increase \(n_0\) to keep these probabilities
uniformly interior.  Let \(q_L\) be the iid Bernoulli \(p_+\)
probability on the actual slots of the left interval that
\[
 cv-aN([s_d,s_d+v])+bN_B([s_d,s_d+v])\ge-C_2U_d,
       \qquad 0\le v\le L_d.
\]
Let \(q_R\) be the iid Bernoulli \(p_-\) probability on the right
interval, read backwards, that
\[
 -cv+aN([t_d-v,t_d])-bN_B([t_d-v,t_d])\ge-C_2V_d.
\]
Define \(q_d=\mathbf1_{\eqref{ld:endpoint-occupancy}}q_Lq_R\).
These iid probabilities are upper bounds for the original conditional
law, not the law of the accepted word.

To prove this, fix all revealed labels and a complete named type.
If \(m\) slots and \(e\) binary successes have been revealed, the
remaining binary word is a uniform \((k-e)\)-subset of \(N=n-m\)
slots.  Its fraction \(p'=(k-e)/(n-m)\) satisfies
\[
 |k-np_0|\le C(X+Y),\qquad
 |p'-p_0|\le C(X+Y+1+d)/n,
\]
so \(p_-\le p'\le p_+\).  Let \(M_J\le N/2\) be the combined
number of slots in both free intervals.  Relative to iid Bernoulli
\(p'\), their joint density under the exact subset law is
\begin{equation}
 \frac{\mathbb P\{\operatorname{Bin}(N-M_J,p')=Np'-j\}}
 {\mathbb P\{\operatorname{Bin}(N,p')=Np'\}}\le C.
 \label{ld:joint-subset-density}
\end{equation}
The numerator is bounded by a maximum atom and the denominator is a
central atom.  This is one comparison for both intervals and holds
for every assignment on them.  The endpoint identity and occupancy
bound show that \eqref{ld:bank-path} requires both displayed tests,
with the stated allowances.  Under the iid law the intervals are
independent.  The left event increases and the right event decreases
with binary successes, so coordinatewise uniform coupling gives
\begin{equation}
 \mathbb P\{\text{accepted path}\mid
              \text{complete type, revealed labels}\}\le Cq_d(G).
 \label{ld:type-bank-upper}
\end{equation}
Additional named restrictions within either binary list do not change
this uniform binary-subset marginal.

For several blocks, fix the whole type and all exposed labels before
multiplying \eqref{ld:type-bank-upper}.  Suppose
\(0\le F\le A(K)\) times the necessary path indicators, where
\(A\) is supported on the original central windows and
\(W=\mathbb E_\pi A\).  Applying Lemma~\ref{ld:reveal} only after
the full type sum gives
\begin{equation}
 \mathbb E_\pi[F\mid Q]\le CW R_{\rm strip}
 \prod_{j\in I}(1+n_j)^{(a_j-1)/2}
 \prod_{\substack{i\notin I\\i\ {\rm selected}}}q_{i,d_i}(G_i),
 \label{ld:class-banks}
\end{equation}
where \(I\) denotes fully exposed blocks and
\[
 R_{\rm strip}=
 \prod_{i\notin I:\,m_i>n_i/2}(1+m_i)^{(a_i-1)/2}.
\]
At active depths, \(m_i\le\varepsilon n_i\), so this extra factor
only concerns the discarded large-depth regime.  Conditioning on a
partial-signature event \(Q\) last averages the uniform bound over
its exact revealed assignments.  There is no factor \(1/\pi(Q)\),
and the type acceptance need not factor across blocks.

\subsection{Moments of the free-interval costs}

For \(d\le D_n\), the functions just constructed satisfy
\begin{equation}
 \mathbb E_{G\mid n}q_d(G)^\eta\le C(1+d)^C\Psi.
 \label{ld:bank-moment}
\end{equation}
Here and below \(A(K)\) and deterministic parameters must be fixed
before these unread clocks for a clock moment assertion.  The class
bound itself is pointwise after fixing the clocks.

The exact-total uniform clock restricted jointly to the two free
intervals has density, relative to a rate \(n/P\) Poisson clock,
\begin{equation}
 \frac{\mathbb P\{\operatorname{Pois}(n(1-\kappa))=n-N_J\}}
 {\mathbb P\{\operatorname{Pois}(n)=n\}}\le C,
 \qquad \kappa=2L_d/P<1/2.
 \label{ld:joint-clock-density}
\end{equation}
Conditional point locations also agree.  Only after this single joint
comparison are the two reference restrictions independent.  Drop the
occupancy indicator for the upper bound.  The auxiliary risk drifts
have magnitude \(O((X+Y+1+d)/n)\), whose product with \(\sqrt{L_d}\)
is bounded.  Theorem~\ref{pr:physical-one}, for either sign,
therefore gives
\begin{equation}
 \mathbb E q_d^\eta\le C
 \min\left\{1,\left(\frac{U_d}{\sqrt{1+L_d}}\right)^{2\chi_+}\right\}
 \min\left\{1,\left(\frac{V_d}{\sqrt{1+L_d}}\right)^{2\chi_-}\right\}.
 \label{ld:two-one-ends}
\end{equation}
The one-end theorem is applied to the original positive mark pair in
both orientations, with compact drift changes; all original spacings
were retained before \eqref{ld:joint-clock-density}.

If \(M_d\le\sqrt n\), then \(1+L_d\asymp n\) and the product
of the reserve ratios is \(\asymp U_dV_d/n\).  If
\(M_d>\sqrt n\), then \(1+L_d\asymp(n/M_d)^2\); the larger
reserve factor is bounded by one and the smaller ratio is
\(\asymp U_dV_d/n\).  The exponent changes satisfy
\[
 |\chi_+-\chi(\eta,\rho_0)|+
 |\chi_--\chi(\eta,\rho_0)|
       \le C\{\mathcal T+(1+d)/n\}.
\]
Every nontrivial reserve ratio in \eqref{ld:two-one-ends} is at least
\(c\min(1,\mathcal T)\).  Multiplying its negative logarithm by
the last bound is uniformly bounded, because
\(\mathcal T|\log\mathcal T|\) is bounded and \(d\le D_n\).
Replace both exponents by \(\chi(\eta,\rho_0)\) and use
\(U_dV_d\le(1+d)^2(1+X)(1+Y)\) to prove
\eqref{ld:bank-moment}.  If the free interval was too short for the
one-end theorem, then \(M_d\gg n\), and \(d=o(\sqrt n)\)
forces one of \(X,Y\) to have order \(n\).  Thus \(\mathcal T\)
has a fixed positive lower bound and the convention \(q_d=1\) is
valid.  Bounded \(n<n_0\) is treated in the same upper estimate;
this does not enlarge the support convention of the source lower.

Fix \(b_0>0\) and an integer \(B_0\) large enough to sum the
later pivot multiplicities.  Choose \(D_0\) with
\(D_0\eta_->2\sup\chi+1\), and set
\begin{equation}
 D^{\rm raw}(G)=(2+n)^{-D_0}
       +\sum_{d\ge0}(1+d)^{B_0}e^{-b_0d^\gamma}q_d(G),
 \qquad D(G)=\min(1,D^{\rm raw}(G)).
 \label{ld:cost-envelope}
\end{equation}
Choose \(C_{\rm cut}\) after \(D_0\) so that the \(\eta\)-powers
of the numerical coefficient tail beyond \(D_n\) sum to at most
\(C(2+n)^{-D_0\eta}\), uniformly in \(\eta\).
Subadditivity and \eqref{ld:bank-moment} then give
\begin{equation}
 (2+n)^{-D_0}\le D\le1,\qquad
 \mathbb E D^\eta\le C\Psi.
 \label{ld:cost-moment}
\end{equation}
At depth zero \eqref{ld:type-bank-upper} gives \(q\le CD\) for
each supported complete type.  Clipping at one preserves this bound
because \(q\le1\).

Large depths are paid using their actual geometric penalty, not a
false small moment for \(q_d=1\).  If that penalty includes
\(cd^\gamma\), increase \(C_{\rm cut}\) so that
\begin{equation}
 e^{-cd^\gamma/2}(1+n)^{H/2}\le C(2+n)^{-D_0},
       \qquad d>D_n.
 \label{ld:large-depth}
\end{equation}
This pays both the selected bridge floor and any large-strip reveal
ratio in that block, leaving half the penalty to sum pivots.  An
unselected block with \(m>n/2\) instead has
\(n\le2m\le C(1+d)\), so its reveal ratio is a polynomial in
its actual depth.  For active depths use
\(q_d\le CD(1+d)^{-B_0}e^{b_0d^\gamma}\), choosing \(b_0\)
below the available geometric coefficient.  When several blocks
share a depth parameter, divide its fixed penalty among the
dimension-bounded number of factors before choosing these constants.

A related count extraction will sometimes be useful in a completely
free interior interval.  Conditional on every label outside it, a
band of at most \(C(1+R)\) possible counts there has probability at
most \(\min(1,C(1+R)/\sqrt{1+N_I})\), irrespective of its center.
This is the binomial maximum-atom bound in
\eqref{ld:multinomial-atoms}.  The whole band is summed before taking
a power.  If this interval and the two path intervals are disjoint
with total length less than half the original block, the single
density comparison \eqref{ld:joint-clock-density} applies to all
three at once.  The count factor has moment
\(O(\min(1,(1+R)/\sqrt{1+P})^\eta)\); atypically small occupancy
has exponentially small probability.  This explains the ordinary
count-window factor without replacing the original class by three
independently conditioned experiments.

\subsection{Named regularity on each retained word}

\begin{lemma}[Uniform named refinement]
\label{ld:named-REG}
Let a fixed parent list of \(n\le C(1+L)\) points in an original
block satisfy, at the fixed grid boundaries,
\begin{equation}
 |N((a,b])-(n/L)(b-a)|\le
 \theta_{\rm in}(b-a)+K_{\rm in}
       \{1+d_L(a)^\zeta+d_L(b)^\zeta\},
 \label{ld:parent-REG}
\end{equation}
where \(0<\zeta<1/2\) and the grid has bounded boundary density.
For every feasible full refinement type \((k_c)\), its original
uniform multiset refinement satisfies, simultaneously in the finitely
many names,
\begin{equation}
 |N_c((a,b])-(k_c/L)(b-a)|\le
 \theta_{\rm out}(b-a)+K_{\rm out}
       \{1+d_L(a)^\zeta+d_L(b)^\zeta\}
 \label{ld:child-REG}
\end{equation}
with conditional probability at least \(1-\varepsilon\).
Here \(\theta_{\rm in}\) is sufficiently small in terms of the
prescribed \(\theta_{\rm out}>0\), and \(K_{\rm out}\) is sufficiently
large in terms of the fixed data and \(\varepsilon\).  The assertion
holds for every fixed parent word, however it was previously selected.
\end{lemma}

\begin{proof}
If an interval contains \(m\) parent positions, its count of name
\(c\) is hypergeometric with mean \(\alpha m\), where
\(\alpha=k_c/n\).  Sampling without replacement has exponential
moments bounded by the corresponding independent Bernoulli moments.
One elementary verification applies the elementary-symmetric-mean
inequality to the \(n\) numbers \(e^t\) or \(1\), and handles
the lower tail by complementation.  Chernoff's argument gives
\[
 \mathbb P\{|N_c-\alpha m|>u\}\le
       2\exp\{-cu^2/(m+u)\}.
\]
The displacement of the mean from \((k_c/L)(b-a)\) is bounded by
\eqref{ld:parent-REG}, and
\(m\le C(b-a)+K_{\rm in}(1+d_L(a)^\zeta+d_L(b)^\zeta)\).
After the stated choices, the probability of violating
\eqref{ld:child-REG} is at most
\[
 C\exp\{-c_1(b-a)-c_2K_{\rm out}
                 (1+d_L(a)^\zeta+d_L(b)^\zeta)\}.
\]
The sum over all grid endpoint pairs is uniformly finite.  Pairs
near an original endpoint are summed by the convergent series
\(\sum_{j\ge0}e^{-cj^\zeta}\); pairs in opposite halves also pay
\(e^{-c_1(b-a)}\).  Increasing \(K_{\rm out}\) makes the finite
union over names and blocks have probability at most
\(\varepsilon\).  Bounded blocks have bounded counts, absorbed
by the fixed intercept; empty lists have their unique empty refinement.
\end{proof}

Apply this lemma to both actual lists of a retained binary word,
conditioning on that word and its complete type.  It supplies named
REG with constant conditional success before any rare bridge is
averaged.  If natural rather than empirical rates are used, a central
type changes the rate by \(O(L^{-1/2})\), absorbed by a spare
linear slope above a fixed length threshold.  The lemma concerns
this per-word refinement only; any later function of the selected
spatial lists remains in its own conditional measure.

\subsection{Admission of the macroscopic recoloring intervals}

Consider a prescribed original block with \(P,n\asymp T\), complete
central fine type \(|K_c-np_c|\le B\sqrt{1+n}\), and a fixed
interior interval \(I\) of length at least \(\rho P\), a fixed
positive distance from both endpoints.  For two specified actual
fine names \(s,t\), require
\begin{equation}
 m=N(I)\ge c_In,\qquad
 |K_{I,s}-p_sm|+|K_{I,t}-p_tm|
                  \le\frac{m}{\log(2+T)}.
 \label{ld:broad-bank}
\end{equation}
The same requirement may be imposed in finitely many prescribed
intervals.  Conditional on \(n\), the uncolored occupancy failure
has probability \(O(e^{-cn})\).  Conditional on the complete type
and a clock with \(m\ge c_In\), the named occupancy is
hypergeometric, with mean \(mK_c/n\).  Its distance from \(p_cm\)
is \(O(m/\sqrt n)\), which is less than
\(m/(8\log(2+T))\) above a fixed threshold.  The same Chernoff
bound as in Lemma~\ref{ld:named-REG} gives, at each full central type,
\begin{equation}
 \mathbb P_{G,\pi}\{\eqref{ld:broad-bank}\text{ fails}\mid K\}
            \le C\exp\{-cT/\log^2(2+T)\}.
 \label{ld:broad-failure}
\end{equation}

Let \(q(K,G)\) already include the allocated curved path, binary REG,
and full named REG, with
\begin{equation}
 q\le CD,\qquad \mathbb E q^\eta\ge c\Psi,\qquad
 \mathbb E D^\eta\le C\Psi,\qquad
 \Psi\ge c_0(2+n)^{-B_1}.
 \label{ld:local-inputs}
\end{equation}
The all-reserve and contracted-chord bridge theorems
\ref{bridge:all-reserves} and \ref{bridge:contracted}, followed by
Lemma~\ref{ld:named-REG}, supply the lower assertion.  The free-interval
construction supplies the other two assertions.  Let \(q_{\rm good}\)
also impose \eqref{ld:broad-bank}, and put
\(q_{\rm bad}=q-q_{\rm good}\).  Subadditivity, then Jensen, give
\begin{align}
 \mathbb E q_{\rm good}^\eta
 &\ge\mathbb E q^\eta-\mathbb E q_{\rm bad}^\eta\notag\\
 &\ge c\Psi-(\mathbb E q_{\rm bad})^\eta
 \ge c\Psi-C\exp\{-c\eta T/\log^2(2+T)\}
 \ge c'\Psi.
 \label{ld:local-bank-retention}
\end{align}
The final inequality follows after one fixed threshold enlargement,
uniformly in the compact range of \(\eta\), from the polynomial
lower bound.  All walls and REG restrictions were already inside
\(q\), and the dominating cost \(D\) does not change.  The same
argument applies to an unselected regular block with \(D=\Psi=1\).
This local fractional estimate is taken before any later scale or
root probability is integrated.
For large \(T\), \eqref{ld:broad-bank} also guarantees at least
\(cn\) occurrences of each of its two names, enough for every
prescribed \(O(\log T)\) recoloring.

\subsection{One coupled count measure, with the actual owner powers}

The count construction in Appendix~\ref{app:endpoint-details} precedes
the unread positions.  More precisely, the diagonal exit bands and
their endpoint count transports (\ref{paper:ep-diagonal},
\ref{paper:ep-count-box}, and \ref{paper:ep-count-transport}) use
only the proper owners.  Projecting out terminal private names and
then filling their prefix lists under the original conditional law
gives the whole prefix count measure of
\ref{paper:ep-prefix-source}.  Conditional on every original prefix
total in its own central window, this is one positive measure \(\mu\)
on full fine types and count-only state, independent of newly unread
prefix positions, with
\begin{equation}
 W=\mu(1)\asymp\prod_{j<m}w_j.
 \label{ld:whole-count-mass}
\end{equation}
It retains the displaced good-name box.  The conditional terminal
completion in \ref{paper:ep-terminal-admission} has constant mass
for each retained prefix type and each suffix root vector in the
specified central construction.  Independence from suffix uncolored
counts follows from the preceding support of the proper-owner
construction.

For each fixed complete type, take the legal pair of
\ref{paper:ep-local-pair}, namely \eqref{lower:local-pair}.  All its
outside contexts have the same actual increment, so it is below
both actual intrinsic endpoint minima and belongs to their common
positive boxes.  Write its path as \(\ell_i^-+E_i\), with
\(\ell_i^-\) the joining chord.  Impose on that original clock
\begin{equation}
 \theta\ell_i^-(t)+E_i(t)\ge
         2g\min(t,L_i-t)^\gamma-h_W,
       \qquad 0<\zeta<\gamma<1/2,\quad0<\theta<1,
 \label{ld:source-wall}
\end{equation}
together with binary REG.  Conditional on this word, refine both
lists by Lemma~\ref{ld:named-REG}.  On unselected blocks require
only the appropriate ordinary and named REG.  Each local mass
\(q_i(\omega,G_i)\) then satisfies \eqref{ld:local-inputs} at
\(\eta_i=\sigma_{\operatorname{owner}(i)}\), with the factor
\(\Psi_i^I\) of \eqref{lower:conditional-factor}.
The necessary flat path follows from \eqref{ld:source-wall} because
its uncontracted chord is larger and nonnegative, so the depth-zero
bank bound applies at these intrinsic scales.  Bounded supported
experiments retain their fixed positive conditional mass and have
\(D_i=\Psi_i^I=1\); a scalar parameter is not assigned to a zero
count.  Applying \eqref{ld:local-bank-retention} incorporates all
prescribed macroscopic recoloring intervals without changing these
estimates.

The following measure identity retains the individual owner powers.
\begin{lemma}[Whole-type mixed moment]\label{ld:mixed}
Let \(\mu\) be any finite positive measure on complete types,
independent of the original independent environments \(G_i\), and
suppose \eqref{ld:local-inputs} holds uniformly in its support, with
possibly different \(\eta_i\in[\eta_-,\eta_+]\subset(0,1)\).
Then, for
\(f(G)=\int\prod_iq_i(\omega,G_i)\,d\mu(\omega)\),
\begin{equation}
 \mathbb E\left[f(G)\prod_iD_i(G_i)^{\eta_i-1}\right]
          \asymp\mu(1)\prod_i\Psi_i.
 \label{ld:mixed-result}
\end{equation}
\end{lemma}
\begin{proof}
For each fixed \(\omega\) and \(i\), H\"older applied to
\[
 q_i^{\eta_i}
  =(q_iD_i^{\eta_i-1})^{\eta_i}
                (D_i^{\eta_i})^{1-\eta_i}
\]
gives
\[
 \mathbb E[q_iD_i^{\eta_i-1}]
 \ge\frac{(\mathbb E q_i^{\eta_i})^{1/\eta_i}}
            {(\mathbb E D_i^{\eta_i})^{(1-\eta_i)/\eta_i}}
 \gg\Psi_i.
\]
The upper bound follows from \(q_i\le CD_i\).  These expectations
are finite by that same upper bound.  At a fixed complete type the
original environments are independent, so their local integrands
can be multiplied.  Tonelli then integrates the resulting uniform
two-sided bounds over the entire \(\mu\).
\end{proof}

For the present prefix,
\(f_{\rm pre}(G)=\int\prod_iq_i(\omega,G_i)\,d\mu(\omega)\).
The owner costs in \eqref{lower:owner-cost} give the deterministic
width multiplier \(\prod_{j<m}w_j^{\sigma_j-1}\).
Multiplying \eqref{ld:whole-count-mass} by this multiplier yields
exactly \(\prod_{j<m}w_j^{\sigma_j}\), while every local path
retains its own owner power.

\subsection{The ordinary suffix and its one count atom}

We detail why the suffix is compatible with the whole prefix
construction and with a later common uncolored target.  Suppose first
that it contains an ordinary weak block.  It is the first suffix
block.  Let \(M_0\) be its length, \(P\) the entire preceding
original length, and \(B\) the following positive-drift length.
The clipped-reference and KKT surplus estimates give
\begin{equation}
 A_{\rm mean}=\sum_{g<i}L_gg_g(r_g)\asymp P,\quad
 B_{\rm mean}=-\sum_{g>i}L_gg_g(r_g)\asymp B,\quad
 A_{\rm mean}+M_0g_i(r_i)=B_{\rm mean}.
 \label{ld:ordinary-balance}
\end{equation}
Thus \(D=A_{\rm mean}\asymp P\),
\(P\le C(M_0+B)\), and \(B\le C(P+M_0)\), including zero-side
cases.  This sums all original old groups and does not assign a
common drift sign to each of them.

At the weak block's left endpoint, reference clipping gives cost
at least \(c\min(P,M_0+B)\gg P\).  The full-prefix context has
cost exactly \(A_{\rm mean}\) and no proper good-group term.
At the right endpoint the same argument gives a lower bound
\(c\min(P+M_0,B)\gg B\), attained up to constants by the
full-prefix context of cost \(B_{\rm mean}\).  Therefore the
intrinsic endpoint minima are
\begin{equation}
 X_0^I\asymp1+P,\qquad Y_0^I\asymp1+B.
 \label{ld:ordinary-small-boxes}
\end{equation}
In both comparisons the common terminal target is initially zero;
there is no incoming-exit reserve in the small endpoint.
The constant-mark theorem \ref{bridge:ordinary-strong}, applied
to the legal representative with actual increment
\(r_iM_0-n\log(r_i+1)\), gives the local scale
\begin{equation}
 \beta_{min}\asymp
     \min\{1,(1+P)(1+B)/(1+M_0)\}.
 \label{ld:weak-beta}
\end{equation}
This uses the deterministic-mark theorem directly.

If the weak block is a largest suffix block, this is the bracket in
\eqref{lower:ordinary-factor}, by \eqref{ld:ordinary-balance}.
If a later positive block is largest, then \(B\ge M_0\), so
\eqref{ld:weak-beta} is bounded below.  Choose the weak drift
threshold smaller than half the fixed positive tail gap.  The mean
of that largest positive block then dominates the possible negative
mean of the weak block, giving \(D\gg M\) and saturating the
bracket in \eqref{lower:ordinary-factor} as well.  If there is no
weak block, every suffix drift is uniformly positive, and
\(D\gg U\); both statements are again saturated.

Condition on a complete retained prefix.  Its uncolored discrepancy
\(\Delta_{\rm pre}\) satisfies
\(|\Delta_{\rm pre}|\le C\sqrt{1+V_{\rm pre}}\), and the Euler
identity gives the suffix target
\begin{equation}
 E_x=D+(x-\Delta_{\rm pre})/r_m,\qquad
 |E_x-D|\le C_A\sqrt{1+D}
 \label{ld:root-target}
\end{equation}
for every prefix-uncolored
\(0\le x\le A\log(2+V_{\rm pre})\), using
\(D\gg V_{\rm pre}\).  Fix the nonreservoir suffix counts in
their own central windows.  In a predetermined largest block use
the actual new count
\begin{equation}
 n_*^x=\left\lfloor
 \frac{\sum_{\rm suffix}c_iL_i+E_x-
                \sum_{i\ne *}w_in_i}{w_*}\right\rfloor,
 \qquad w_i=\log(r_i+1),\quad c_i=r_i.
 \label{ld:new-root}
\end{equation}
The support threshold and central root construction of
Theorem~\ref{bridge:ordinary-root} make it nonnegative and central,
with Poisson mass \(\asymp(1+U)^{-1/2}\).  All its original
\(n_*^x+1\) gaps are then sampled.  In the unconditional Poisson
experiment this means retaining the event that the sampled total
equals \eqref{ld:new-root}; it does not alter a previously fixed
tuple.  The bounded floor error is \(-C\le\epsilon\le0\).

The count \(n_*^0\) obtained with \(x=0\) is used only to verify
the zero-target prefix boxes through the uniform count completion
in \ref{paper:ep-terminal-admission}.  No clock at that virtual
count is sampled or reused.  Increasing \(x\) preserves the
prefix local endpoint lower bounds: terminal early bounds gain
the common entrance, and terminal late intrinsic bounds have
neither entrance nor incoming carry.  Thus prefix paths, costs,
and local pairs are fixed before \(x\) is evaluated.

If the weak block is not the reservoir, its local count is unchanged.
If it is, put \(k=n_*^x-n_*^0\ge0\).  Then \(k\le C(1+x)\),
and the max formula for the legal pair gives
\begin{equation}
 \Delta_x=\Delta_0-w_*k,\quad
 x_0^-\le x_x^-\le x_0^-+w_*k,\quad
 a_*Y_0^I\le y_x^-\le y_0^-.
 \label{ld:root-pair-change}
\end{equation}
Equation \eqref{ld:ordinary-small-boxes} absorbs
\(w_*k\le C_A\log(2+P)\) in the left box, while the right
endpoint keeps its positive lower box.  Its scale is still
\eqref{ld:weak-beta}; an incoming exit or \(x\) has not been
added to its small right endpoint.

Complete the terminal full type at its constant conditional mass,
then impose the weak ordinary representative's allocated strong wall
and uncolored REG.  The other suffix blocks have separated drift.
Reading a positive slope from the left, or a negative slope from
the right, their REG implies the allocated inward wall: absorb
\(K(1+t^\zeta)+gt^\gamma\) into a fixed intercept and a fraction
of the linear drift.  Long central blocks preserve the drift gap;
bounded supported blocks use their fixed finite convention.  The
same argument pays the strong and reference-paid native comparisons.
After counts are fixed, the original suffix clocks are independent;
there is at most one weak ordinary probability.  Fill named REG
on each retained clock by Lemma~\ref{ld:named-REG}.  This constructs
a joint completion of order
\(\mathcal B_I=(1+U)^{-1/2}\beta_0\) on the same word.

For the matching upper and the support needed later, retain the
lawful full-stop context for every suffix \(q=0\) comparison:
all eligible names remain until its moving cut, then disappear.
Its intrinsic cost is the original reverse suffix risk, with no
incoming reserve.  At a fixed count its centered fluctuation agrees
with that of the minimum representative; the latter's positive
chord, including its \(\theta\)-contraction, lies below the
full-stop chord.  Thus the imposed weak and paid strong walls
imply on their actual joint support
\begin{equation}
 h_{\rm base}+R_{\rm tail}(v)\ge0\qquad(0\le v\le U)
 \label{ld:global-ordinary-wall}
\end{equation}
for one fixed reserve, large enough for floor and fixed intercept
errors.  Drop other filters and apply the all-count upper in
Theorem~\ref{bridge:ordinary-root} at this reserve and target
\eqref{ld:root-target}.  It gives \(O(\mathcal B_I)\), with the
reservoir denominator retained before its Poisson supremum.
Consequently, for every retained complete prefix word and every
permitted \(x\), its actual suffix completion kernel is comparable
to \(\mathcal B_I\).  Integrating the private prefix gives the
pointwise relation
\begin{equation}
 \mathbb E_{\rm suffix}f_0^x(G_{\rm pre},G_{\rm suffix})
        \asymp\mathcal B_I f_{\rm pre}(G_{\rm pre}).
 \label{ld:suffix-kernel}
\end{equation}
This expectation includes the one root atom, all suffix positions,
and the original named refinements.  It may be multiplied by any
nonnegative prefix-uncolored cost, but it does not scalarize a
future private-list weight.

Apply \eqref{ld:suffix-kernel} first, then Lemma~\ref{ld:mixed}
at each full type, and finally \eqref{ld:whole-count-mass}.  This
proves the source mixed moment
\begin{equation}
 \mathbb E\left[f_0^x\prod_j(D_j^{\rm own})^{\sigma_j-1}\right]
 \asymp\mathcal B_I\prod_{j<m}w_j^{\sigma_j}
                  \prod_{i\in\mathcal I}\Psi_i^I(\sigma_{\rm owner(i)}).
 \label{ld:source-mixed}
\end{equation}
The ordinary factor is linear and occurs once.

\subsection{Targets fixed by the costs and a literal fine-word density}

We specify the order of constants before defining the adaptive targets.
First fix compact scalar ranges, \(\theta\), \(0<\zeta<\gamma<1/2\),
and the finite count and displacement support.  Choose a polynomial
floor exponent \(D_0\) as in \eqref{ld:cost-moment}; it depends on
the compact exponent range and number of factors, not on later fixed
wall constants.  Use the literal widths \(w_j\le1\), with
\(-\log w_j\le\tfrac12\log T_j\).  The floors imply predetermined
bounds
\begin{equation}
 y_j:=-\log D_j^{\rm own}\le C_0+C_1\log(2+T_j),\qquad
 y_j\ge0.
 \label{ld:targets}
\end{equation}
Only selected head blocks enter the terminal cost, so its target
is bounded in terms of \(V_{\rm pre}\).  The hierarchy
\(T_l\le C_HT_j\) for \(l<j\) and bounded reciprocal child
entropies now fix all logarithmic transfer ranges described below.
Construct the common upper costs simultaneously for these full ranges,
then choose the source with spare strong-curve and REG margins.
Finally enlarge the macro-block and root support thresholds.  These
later constants do not alter the floor exponent or logarithmic target
bounds.

Set
\begin{equation}
 x=\sum_{j=1}^m r_jy_j,\qquad r_j=\sigma_j/\sigma_1.
 \label{ld:common-target}
\end{equation}
This is prefix-uncolored and lies in the preassigned range of
\eqref{ld:root-target}.  Choose the actual root by
\eqref{ld:new-root}, obtaining \(X_1=x+\epsilon\),
\(-C\le\epsilon\le0\).  Let \(F_0^x\in[0,1]\) be the joint
source acceptance just constructed, with private mass \(f_0^x\).
Proper targets are still zero at this stage.  The suffix is sampled
only once, and \eqref{ld:source-mixed} applies to this \(x\).

Let \(A_j\) be the logarithm of the first shallow effective
alphabet of proper owner \(j\), and put
\(b_j=u_jA_{j+1}\).  In its prescribed macro block take the
immediate child's first shallow fine label \(c_j\), and that
block's eligible shallow exterior fine label \(t_j\).  Both are
actual original labels with probabilities bounded below, and their
ratio is
\[
 p(c_j)/p(t_j)=\exp\{(u_j-1)A_{j+1}\}.
\]
The destination need not be the first name of the entire owner.
Define, using only the uncolored prefix,
\begin{equation}
 e_0=0,\quad
 e_j=\left\lceil\frac{y_j+A_je_{j-1}}{b_j}\right\rceil,\quad
 \delta_j=b_je_j-A_je_{j-1}-y_j\in[0,b_j),\qquad j<m.
 \label{ld:cascade}
\end{equation}
Equation \eqref{ld:targets} and the hierarchy give the predetermined
bound \(e_j\le A'_H\log(2+T_j)\).  Condition
\eqref{ld:broad-bank} and the final threshold ensure that all these
occurrences exist.  Choose a uniform \(e_j\)-subset of the
\(c_j\)-slots in its bank and recolor it \(t_j\).  The banks lie
in disjoint original blocks.  This defines a single probability
kernel on the full fine word, preserving every clock and the root.

Here is its exact density calculation.  In one bank, write \(s,t\)
for source and destination names, and \(k',l'\) for their output
counts.  Every preimage is obtained by choosing \(e\) of the
output \(t\)-slots, and has \(k'+e\) available source slots.
Its original word probability is larger by \((p_s/p_t)^e\).
Thus the unrestricted push-forward density is
\begin{equation}
 R_e(k',l')=(p_s/p_t)^e
       \frac{\binom{l'}e}{\binom{k'+e}e}.
 \label{ld:fine-RN}
\end{equation}
On an output with a retained preimage, expand the \(e\) factors.
The two bank counts differ from \(p_sm,p_tm\) by at most
\(m/\log(2+T)\), and the shift is at most \(e\).  Since both
probabilities are uniformly positive,
\begin{equation}
 |\log R_e|\le C\{e/\log(2+T)+e^2/m\}\le C_A.
 \label{ld:fine-RN-bound}
\end{equation}
Deleting preimages can only decrease the density.  Multiplying the
upper bounds over the finitely many disjoint banks gives \(C_R\).
If \(\nu_y\) is the push-forward of \(\pi F_0^x\), define
the actual accepted target by
\begin{equation}
 F_y^I=\frac1{C_R}\frac{d\nu_y}{d\pi},\qquad
 0\le F_y^I\le1,\qquad
 \mathbb E_\pi F_y^I(G,\cdot)=f_0^x(G)/C_R.
 \label{ld:literal-density}
\end{equation}
The mass identity holds for every fixed full uncolored environment,
including its root indicator and the adaptive vector \(e(G)\).
Thus \eqref{ld:literal-density} preserves the integrated mass.

\subsection{The reconstructed support and its common upper costs}

We verify the necessary target paths directly on the changed slots.
Fix a selected block of length \(P\) that is recolored.  Its
canonical retained binary set contains \(s\) and excludes \(t\).
Every target full type \(K\) has the same inverse source type
\(K-\sum_je_jv_j\), where \(v_j\) is its source-to-destination
count shift.  Thus its source pair \(x_0,y_0\) is recovered
without knowing which slots were changed.  If \(L_{\rm flip}(t)\)
counts recolorings before time \(t\), the reconstructed process
satisfies the exact identity
\begin{align}
 \widetilde Z(t)
  &=x_0+(c+be/P)t-aN(t)+bN_{B,\rm target}(t)\notag\\
  &=Z_{\rm source}(t)-b\{L_{\rm flip}(t)-et/P\},\qquad
 \widetilde Z(0)=x_0,\quad\widetilde Z(P)=y_0.
 \label{ld:slot-identity}
\end{align}
This follows from \(N_{B,\rm target}=N_{B,\rm source}
-L_{\rm flip}\).

For every possible chosen subset in the fixed interior bank,
\begin{equation}
 |L_{\rm flip}(t)-et/P|\le
         CeP^{-\gamma}\min(t,P-t)^\gamma.
 \label{ld:slot-error}
\end{equation}
Before the bank the left side is \(et/P\), after it it is
\(e(P-t)/P\), and inside it the depth is \(\gg P\) and
the left side is at most \(e\).  The spare \(2g\) curve in
\eqref{ld:source-wall} and \(e=O(\log P)\) therefore imply
\begin{equation}
 \theta\ell_0(t)+E_{\rm target}(t)
       \ge g\min(t,P-t)^\gamma-h_W
 \label{ld:transported-wall}
\end{equation}
above the final threshold.  Its endpoints are exactly the original
source pair.  On an unflipped native block the word and local type
are unchanged, and the inverse complete type recovers the same
pair.  A previous owner's recoloring preserves every point of the
later native clock, so the estimate uses that block's own length.

Named REG is preserved with the predetermined spare constants.
Any interval meeting the macro bank either has length \(\gg P\)
or has an endpoint of depth \(\gg P\).  Its \(O(\log P)\)
count change is paid respectively by a spare linear slope or a
spare depth-to-the-\(\zeta\) intercept.  The empirical-center
shift \(O(\log P/P)\) is paid by the same margins.  The terminal
suffix has no recoloring bank; its root, clocks, and fixed-reserve
ordinary wall remain unchanged.  The support proved here is the
source-reconstructed support \eqref{ld:transported-wall}.  The
separate same-word information and carry comparison in
Appendix~\ref{app:history-details} transfers full fine histories.

To construct one bank envelope in advance for all the permitted
transfers, retain \(U_d,V_d,L_d\) of \eqref{ld:bank-length} and
let \(E_{max}=\lfloor A\log(2+T)\rfloor\).  Use
\begin{equation}
 p_\pm=p_0\pm C_1(X+Y+1+d+E_{max})/n,\qquad
 c_{max}=c+bE_{max}/P.
 \label{ld:all-e-parameters}
\end{equation}
The left free interval uses speed \(c_{max}\) and probability
\(p_+\); the reversed right interval uses speed \(c\) and
probability \(p_-\).  These opposite choices both enlarge their
necessary events.  The allowances remain \(C_2U_d,C_2V_d\).
After fixing the true revealed assignments and full target type,
\eqref{ld:slot-identity} gives
\(\lvert K_B-np_0\rvert\le C(X+Y)+e\); thus its remaining
binary fraction lies between \(p_-\) and \(p_+\).
The joint exact-subset comparison \eqref{ld:joint-subset-density}
then applies to both intervals on the original target word.  This
argument uses the direct support \eqref{ld:transported-wall}.

The joint clock comparison \eqref{ld:joint-clock-density} is
unchanged.  The new drift term times \(\sqrt{L_d}\) is at most
\(C\log(2+T)/\sqrt n\), uniformly bounded because the recolored
block is macroscopic.  The additional exponent error is
\(O(E_{max}/n)\), and its product with the logarithm of any
active reserve ratio is \(O(\log^2(2+n)/n)\).  The proof of
\eqref{ld:bank-moment} therefore gives
\begin{equation}
 \mathbb E(q_d^*)^\eta\le C(1+d)^C\Psi_{\rm source},\qquad
 \mathbb E(D^*)^\eta\le C\Psi_{\rm source},
 \label{ld:all-e-moments}
\end{equation}
using the same floor and coefficient tail.  All costs here are fixed
from the entire range \(E_{max}\), before any target is evaluated,
with the original source reserves.

For the class upper, use the count-only prefix acceptance
\(A_{\rm pre}\), extended by one on suffix names, and drop the
source's suffix root and terminal filters.  The full \(F_0^x\)
itself is not a clock-independent count weight.  Every preimage of
a target type has the same inverse type, so positivity and
\eqref{ld:fine-RN-bound} bound \(F_y^I\) by the shifted weight
\[
 A_e(K)=A_{\rm pre}(K-\textstyle\sum_je_jv_j).
\]
Its whole original type mass \(W_e\) satisfies
\begin{equation}
 cW\le W_e\le CW,\qquad W=\mathbb E_\pi A_{\rm pre}.
 \label{ld:shifted-band}
\end{equation}
Indeed, in one changed block the adjacent multinomial ratio is
\[
 \frac{\pi(k+ev)}{\pi(k)}
 =(p_t/p_s)^e\prod_{l=0}^{e-1}\frac{k_s-l}{k_t+l+1},\qquad
 \left|\log\frac{\pi(k+ev)}{\pi(k)}\right|
       \le C(e/\sqrt n+e^2/n).
\]
All shifted types are feasible above the fixed threshold, and
unchanged block factors cancel exactly.  Shift the entire coupled
type sum once to prove \eqref{ld:shifted-band}.  The bound holds
simultaneously in every allowed \(e\); it therefore remains
pointwise for \(e(G)\), without treating \(A_e\) as independent
of \(G\).  Apply the original conditional free-interval argument
at a fixed whole type, sum this whole shifted acceptance, and
replace \(W_e(G)\) by \(CW\) before averaging clocks.  This proves
\eqref{ld:class-banks} for the actual transported density with the
all-transfer envelopes.  Its count factor contains no root atom;
that factor remains in the separate source-mass calculation.

\subsection{Exact exits, likelihood, and the first mass}

The grouped residuals change under \eqref{ld:cascade} by
\begin{equation}
 Z_j^y-Z_j^0=b_je_j-A_je_{j-1}=y_j+\delta_j\quad(j<m),\qquad
 Z_m^y-Z_m^0=-A_me_{m-1}.
 \label{ld:output-changes}
\end{equation}
Their \(r_j\)-weighted sum is zero because
\(r_jb_j=r_{j+1}A_{j+1}\).  Define the new exits by the genuine
identity \(Z_j^y=y_j+\eta_j^y-\eta_{j-1}^y/u_{j-1}\).
Subtract the source identity, solve its triangular recursion, and get
\begin{equation}
 \eta_j^y-\eta_j^0=
       \sum_{l\le j}(\sigma_l/\sigma_j)\delta_l=O_H(1).
 \label{ld:exit-changes}
\end{equation}
All these changes are nonnegative.  The positive source exit bands
remain positive and, since \(d_jR_j\ll1\),
\begin{equation}
 0\le E_y:=\sum_{j<m}d_j\eta_j^y\le C_H,\qquad
 Z_m^y=y_m+\epsilon/r_m-\eta_{m-1}^y/u_{m-1}.
 \label{ld:terminal-carry}
\end{equation}
The latter follows from the unchanged root
\(X_1=\sum_jr_jy_j+\epsilon\).  In particular, its \(y_m\)
comes from the actual root; no terminal private transfer is made.

Here is the exact likelihood calculation, including its deterministic
normalization.  Put
\(B_{\rm geom}=\sum_i(h_i-1)L_i=\sum_{j=1}^kV_j\).
The full recursive fine-word law, including the terminal free labels,
satisfies, term by term,
\begin{equation}
 -\log\pi(\omega)=B_{\rm geom}+\sum_{j=1}^mZ_j(\omega),\qquad
 X_1=\sum_{j=1}^mr_jZ_j.
 \label{ld:word-likelihood}
\end{equation}
At an exterior choice the negative log probability is its local
log-alphabet; at a child route it is that log-alphabet minus the
child's weighted log-alphabet.  The child contributions cancel
against the next stage, proving both identities.  Equal-product
internal margins cancel in grouping, by complementarity in
\eqref{hier:eq:KKT}.

Let \(q_j=\sigma_j/\sigma_m-1\), so \(q_m=0\).  Eliminating the
terminal output in \eqref{ld:word-likelihood} gives
\[
 1=e^{B_{\rm geom}+X_1/r_m}\pi(\omega)
                  e^{-\sum_{j<m}q_jZ_j}.
\]
Since \(q_j-q_{j+1}/u_j=d_j\), the exit identities give
\(\sum_{j<m}q_jZ_j=\sum_{j<m}q_jy_j+E_y\).
Therefore
\begin{align}
 (1/r_m-\sigma_1)X_1-\sum_{j<m}q_jZ_j
   &=\sum_{j=1}^m(1-\sigma_j)y_j-E_y
                 +(1/r_m-\sigma_1)\epsilon.
 \label{ld:likelihood-cancellation}
\end{align}
This uses
\((1/r_m-\sigma_1)r_j-q_j=1-\sigma_j\), including \(j=m\).
No exponential in the logarithmic target is discarded.
Also, for two retained words in the same uncolored environment,
\eqref{ld:word-likelihood} and the sum of their exit identities give
\begin{equation}
 \frac{\pi(\omega)}{\pi(\omega')}
             =\exp\{E_y(\omega)-E_y(\omega')\}.
 \label{ld:relative-likelihood}
\end{equation}
It is bounded by \eqref{ld:terminal-carry}, while its actual class
mass still requires the separate class estimates.

For a complete cell occupancy \(g\), put
\(W_g=\prod_{i,c}|c|^{g_{i,c}}/g_{i,c}!\) and
\(Z_F(g)=\sum_\omega F_y^I(g,\omega)\).  Under the natural
Poisson law the cell intensity is \(Z_i=e^{G_i}\).  Expanding
that law in
\(e^{-\sum_i h_iL_i}\sum_gW_gZ_F(g)\) gives deterministic factor
\(e^{-\sum_i h_iL_i+\sum_i Z_iL_i}\), random factor
\(e^{-\sum_iG_in_i}\), and \(\pi(\omega)^{-1}\).
Euler's identity gives
\(\sum_iG_in_i=\sigma_1X_1+\sum_js_jV_j\).
Combining with \eqref{ld:word-likelihood}, the deterministic exponent
is exactly \(-J(L;s)\).  Equation
\eqref{ld:likelihood-cancellation} proves
\begin{align}
 e^{-\sum_i h_iL_i}\sum_gW_gZ_F(g)
   =e^{-J(L;s)}\mathbb E_{\rm nat}\mathbb E_\pi
       \left[F_y^I e^{-E_y}
       \prod_j(D_j^{\rm own})^{\sigma_j-1}
       e^{(1/r_m-\sigma_1)\epsilon}\right].
 \label{ld:exact-first-mass}
\end{align}
This is algebra on the actual count/root support.  The factors
involving \(E_y\) and \(\epsilon\) are bounded above and below.
Using the pointwise mass identity \eqref{ld:literal-density}
inside this same expectation, followed by
\eqref{ld:source-mixed}, evaluates it as
\begin{equation}
 e^{-J(L;s)}\mathcal B_I\prod_{j<m}w_j^{\sigma_j}
             \prod_{i\in\mathcal I}\Psi_i^I(\sigma_{\rm owner(i)}).
 \label{ld:conditional-first-mass}
\end{equation}
The root indicator occurs once, linearly.

To integrate the remaining original prefix totals, let
\(\mu_i=Z_iL_i\).  On their own central windows,
\(\lvert n_i-\mu_i\rvert\le B\sqrt{1+\mu_i}\).
The deterministic intrinsic score minima do not depend on \(n_i\).
Replacing \(n_i\) by \(\mu_i\) in
\eqref{lower:conditional-factor}, with products frozen, changes
\(p_i^0,\rho_i^0\) and the logarithm of its base by
\(O(\mu_i^{-1/2})\) on the compact rare branch.  The base is
at least \(c/(1+\mu_i)\), so local Lipschitz continuity of
\(\chi\) gives
\begin{equation}
 |\log\Psi_i^I-\log\overline\Psi_i^I|
              \le C\log(2+\mu_i)/\sqrt{\mu_i}\le C.
 \label{ld:count-comparison}
\end{equation}
If a branch threshold is crossed, both bases are comparable to that
fixed threshold and the same conclusion holds.  The dimension-bounded
number of original Poisson central windows has constant joint mass;
the source vanishes off them.  Integrating
\eqref{ld:conditional-first-mass} gives \eqref{lower:nominal-mass}.
Fixed positive allowances can then be changed to the profile's normalization because all scores
are finite minima of sums of nonnegative terms, and all exponents
remain in compact ranges.

\subsection{Construction on complete cells}
\label{ld:cells}

We finish by constructing the same measure directly on the complete
cell grid used for arithmetic realization.  This point is needed
because taking a conditional expectation of an old continuous
acceptance would not preserve its fractional moments or its
comparison-class estimates.

Before counts, positions, or labels are sampled, fix a partition of
each original block into cells with endpoints including its physical
endpoints, bounded mesh \(\Delta_*\), and bounded boundary density.
A cell word records the occupancy \(g_{i,c}\) and its exchangeable
labeled slots.  Every final event is invariant under permutations
within a cell.  Only the declared bounded experiments are used for
short blocks, with bounded total, feasible types, and fixed positive
retained mass; these blocks are not recoloring banks.  A later
numerical cutoff for a bank upper estimate does not expand this
bounded support convention.

First consider any lift of a fixed cell word to distinct positions
inside its cells.  For positive marks and either orientation,
\(Z(t)=x+\varepsilon(\sum_{s\le t}W_s-ct)\), monotonicity of
jumps gives, on \([l,r]\),
\begin{equation}
 Z(t)\ge Z(l)-c\Delta_*\quad(\varepsilon=1),\qquad
 Z(t)\ge Z(r)-c\Delta_*\quad(\varepsilon=-1).
 \label{ld:boundary-wall}
\end{equation}
Use the same half-open convention for all boundary counts.  There
is no dependence here on the number of points in the cell.  Since
\(g\min(t,L-t)^\gamma\) oscillates by at most
\(g\Delta_*^\gamma\), a boundary strong wall implies the
continuous strong wall on every lift after increasing its allowance
by \(c_{max}\Delta_*+g\Delta_*^\gamma\).  The reverse inclusion
requires no loss.  Logarithmic soft walls are handled by their fixed
local modulus of continuity.

For an allocated bridge \(\theta\ell+E\), its endpoints are
\(\theta x,\theta y\), and its signed-jump representation has speed
\begin{equation}
 c_\theta=\theta c+(1-\theta)S_{\rm total}/L,
 \label{ld:allocated-speed}
\end{equation}
where \(S_{\rm total}\) is its positive mark sum.  On the large
central experiments this speed stays in a fixed positive compact
interval, independent of the sizes of the smaller intrinsic endpoints.
The same is true after the compensating change \(c+be/L\).
Thus \eqref{ld:boundary-wall} also applies to the allocated walls
at these smaller pairs.

Boundary REG implies continuous REG on every lift.  To bound the
count on \((s,t]\), expand it to its neighboring outer boundaries
for the upper bound and shrink it to the inner boundaries for the
lower.  Length and mean change by at most \(2\Delta_*\) and
\(2\lambda_{max}\Delta_*\).  Also
\[
 |d_L(s)^\zeta-d_L(s')^\zeta|\le\Delta_*^\zeta
\]
for a neighboring boundary \(s'\).  If the inner interval is
empty, its length is at most \(2\Delta_*\), and nonnegativity
gives the required lower bound.  This proves the assertion with
fixed additional intercept, even though an interior cell's permitted
occupancy can grow with its depth.

Define the cell source using boundary allocated walls, boundary
binary REG, and its count box from the outset.  The continuous
relative strong/REG event at the same exact pair is contained in
this boundary event.  The reverse inclusion just proved provides
the continuous necessary upper wall with a fixed larger allowance.
Changing such a fixed allowance changes the all-reserve factors by
fixed constants only.  For named filling, condition on a complete
binary word and refine both actual lists.  Lemma~\ref{ld:named-REG}
applies to each lift with the enlarged fixed REG constants.  A
boundary binary word is a mixture of these lifts; integrating their
uniform success bound gives the same relative retention for the
named boundary event.

Construct fresh costs as cell functions.  At each active depth choose
all strips and free-interval endpoints at original boundaries and
round \(L_d\) down by at most two meshes.  A length below the
fixed threshold uses the already justified \(q_d=1\) convention.
Use the all-transfer parameters \eqref{ld:all-e-parameters} and
test its two necessary walls at boundaries, with only the fixed
mesh allowance.  Denote the resulting product by
\(q_{i,d}^{\rm cell}(g_i)\).  Both the exact-subset comparison
\eqref{ld:joint-subset-density} and the exact-clock comparison
\eqref{ld:joint-clock-density} push to cells with the same density
bound.  The either-sign pointwise implication
\eqref{ld:boundary-wall} then gives
\[
 \mathbb E_{g_i\mid n_i}(q_{i,d}^{\rm cell})^\eta
       \le C(1+d)^C\Psi_i^I(\eta),\qquad d\le D_{n_i}.
\]
Use these functions in \eqref{ld:cost-envelope}, including its
literal floor, numerical tail, and clipping, to define
\(D_i^{I,\rm cell}(g_i)\).  The proofs of
\eqref{ld:cost-moment} and \eqref{ld:all-e-moments} give
\begin{equation}
 q_i^{I,\rm cell}(\omega,g_i)\le CD_i^{I,\rm cell}(g_i),\quad
 \mathbb E(q_i^{I,\rm cell})^\eta\ge c\Psi_i^I(\eta),\quad
 \mathbb E(D_i^{I,\rm cell})^\eta\le C\Psi_i^I(\eta).
 \label{ld:cell-local}
\end{equation}
The middle lower bound follows from the contained continuous event,
with its proved finite scalar support in bounded-count cases.

The macroscopic recoloring intervals are unions of complete cells,
still a fixed positive distance from both block endpoints.  Their
occupancy is binomial conditional on \(n_i\), and the two-name
counts are hypergeometric conditional on the full type and occupancy.
Thus \eqref{ld:broad-failure} and the local subtraction
\eqref{ld:local-bank-retention} apply directly to these cell events.
They preserve \eqref{ld:cell-local} before the coupled type or
root is integrated.

The exact uncolored cell law at fixed total is
\begin{equation}
 \mathbb P(g_i\mid n_i)
       =n_i!\prod_c\frac{(|c|/L_i)^{g_{i,c}}}{g_{i,c}!}.
 \label{ld:exact-cell-law}
\end{equation}
It is the push-forward of the original uniform-point law, equivalently
of its full \((n_i+1)\)-spacing simplex law.  Conditional on the
occupancy, positions are uniform inside cells, labels have their
original categorical law, and conditional on full type the labeled
slots have their uniform multiset law.  By construction,
\[
 q_i^{I,\rm cell}(\omega,G_i)
        =q_i^{I,\rm cell}(\omega,g_i),\qquad
 D_i^{I,\rm cell}(G_i)=D_i^{I,\rm cell}(g_i).
\]
Thus integration over a cell fiber preserves \(q_i\),
\(q_i^\eta\), and \(q_iD_i^{\eta-1}\) exactly.  At a fixed
complete type the block experiments remain independent; Lemma~\ref{ld:mixed} therefore applies before
the full measure \(\mu\) is integrated.

The owner costs, \(y_j\), and \(x\) are now functions of prefix
uncolored cell counts alone.  For every retained prefix cell word,
boundary-to-continuous REG supplies the hypotheses of the suffix
argument on each lift.  The actual new reservoir total is still
\eqref{ld:new-root}.  Sample its cell occupancy directly by
\eqref{ld:exact-cell-law}.  Counts, endpoint identities, and the
bounded root error are unchanged.  Impose the weak ordinary
representative's boundary strong wall and boundary REG, with the
same constant-mass named completion.  The continuous ordinary
lower is contained in it; strong blocks are paid by their separated
drifts.  Conversely the intrinsic full-stop contexts imply the
global ordinary wall at boundaries, and
\eqref{ld:boundary-wall} extends it with one fixed allowance.
The all-count root upper still applies.  Consequently
\begin{equation}
 \mathbb E_{\rm suffix}f_0^{I,\rm cell,x}
       (g_{\rm pre},g_{\rm suffix})
       \asymp\mathcal B_I f_{\rm pre}^{I,\rm cell}(g_{\rm pre})
 \label{ld:cell-root-kernel}
\end{equation}
pointwise on the retained prefix.  Its root is a single count atom,
and the endpoint change of a weak reservoir is still exactly
\eqref{ld:root-pair-change}.

Apply the recoloring kernel to these cell words.  Its uniform subset
choice commutes with every within-cell permutation, as do the source
count, wall, REG, bank, and root events.  Formula \eqref{ld:fine-RN}
and its bound remain exact, so
\begin{equation}
 F_y^{I,\rm cell}=C_R^{-1}d\nu_y/d\pi\in[0,1],\qquad
 \mathbb E_\pi F_y^{I,\rm cell}(g,\cdot)
                         =f_0^{I,\rm cell,x}(g)/C_R
 \label{ld:cell-density}
\end{equation}
is cell-invariant.  At every boundary the same slot identity
\eqref{ld:slot-identity} holds.  Macro interiority and the spare
strong wall prove \eqref{ld:transported-wall} at these boundaries,
then \eqref{ld:boundary-wall} extends it to every lift with a fixed
allowance uniform in the batch size and target.  Named boundary REG
is retained by the same macro-length or macro-depth argument, and the terminal suffix is unchanged.

The output and exit identities
\eqref{ld:output-changes}--\eqref{ld:terminal-carry} are count
identities and incur no cell-rounding error.  Apply
\eqref{ld:cell-root-kernel}, then the local mixed moments from
\eqref{ld:cell-local}, and finally the whole count mass
\eqref{ld:whole-count-mass}.  Equation \eqref{ld:cell-density}
transfers that mass with its fixed normalization.  The exact
likelihood calculation \eqref{ld:exact-first-mass} and the
central-count comparison \eqref{ld:count-comparison} consequently
give the same \(e^{-J}\) times the factors in
\eqref{lower:nominal-mass}.

All constants can be fixed in the order stated before
\eqref{ld:targets}, with the mesh bound included among the initial
data: floor exponent and literal widths first, the logarithmic
transfer ranges next, then free-interval envelopes and spare
strong/REG/refinement and mesh allowances, and the final macro,
total, and root thresholds last.  The subsequent arithmetic
realization may therefore use this bounded cell-invariant density
without reading actual prime values inside its slots.
